\documentclass{zhvt-classic}

% 使用本地源流明体（GenRyuMin2 TW）按文件路径加载，不依赖系统字体名
% 版框（内框／界栏／书口侧栏／鱼尾）改用与版心印章同一印泥棕色：
%   印章图片 yshs.png 实色 = RGB(135,68,52) = HTML 874434
\usepackage{xcolor}
\usepackage{tikz}
% sealbrown 双注册：xcolor 供 TikZ/pgf，L3 供 cls 版框
\definecolor{sealbrown}{HTML}{874434}
\ExplSyntaxOn
\color_set:nnn { sealbrown } { HTML } { 874434 }
\ExplSyntaxOff

\zhvtset{
  grid color       = sealbrown,
  main font        = GenRyuMin2TW-SB,
  main font path   = C:/Users/Administrator/texmf/fonts/opentype/public/GenRyuMin2TW/,
jiazhu font      = GenRyuMin2TW-M,
  jiazhu font path = C:/Users/Administrator/texmf/fonts/opentype/public/GenRyuMin2TW/,
  % 页码框容纳 4 字：正常页「九十七」+前置空 1 格=4 行恰好；
  % 100+ 页「一百零一」共 5 行超出，末字「一/二/三」被框裁掉（2026-10-04 实测），故放宽到 5
  shukou page number chars = 5,
  % 书口书名框高从默认 6 字减到 4 字，把省出的 2 字高拨给 chapter 框，
  % 使 6 字书名+卷名放进 9 字高的 chapter 框后不再压线（2026-10-05）
  shukou title chars = 4,
}

\title{史记}[田叔列传]
\author{汉·司马迁}
\maker{闲叟漫记·重製}


% ===== 索隐述赞样式（2026-11-08 四次定稿 · 愚大人亲手改定参数）=====
% ⚠ 竖排前提：自定义字体家族必须带 RawFeature={vertical:+vert}；
%   simhei.ttf 加 +vhal 会整体加载失败静默回退，只写 +vert；simfang 同理。
\newCJKfontfamily{\shuzanfont}[Path=C:/Windows/Fonts/,RawFeature={vertical:+vert}]{simhei.ttf}
\newCJKfontfamily{\shuzanbody}[Path=C:/Windows/Fonts/,RawFeature={vertical:+vert}]{simfang.ttf}
% 题头：152x30pt 圆角 6pt 棕底白字；baseline (0,2) 使整框落入界行格子不压版框线；
% 白字节点 x 不补偿、y +8.2pt（像素级墨迹实测 dx/dy≈0，L/R 留白约 20pt）。
\newcommand{\shuzanhead}[1]{%
  \begin{tikzpicture}[baseline={(0,2)}]
    \fill[sealbrown, rounded corners=6pt]
      (-76pt,-15pt) rectangle (76pt,15pt);
    \node[text=white, inner sep=0pt, anchor=center] at (0pt,8.2pt)
      {\shuzanfont\fontsize{23}{23}\selectfont #1};
  \end{tikzpicture}%
}
% ★ 修正书口卷名位置：chapter 框起始位置改回 cls 默认值，避免全局影响。
%   只在每个卷首页通过 \juan 命令里的空命令微调该页位置（安全，不会跨页污染）。
\ExplSyntaxOn
\AtBeginDocument{
  \dim_set:Nn \l__zhvt_shukou_chapter_title_start_dim
    { \l__zhvt_shukou_title_height_dim + \l__zhvt_font_size_dim * 2 }
  \dim_set:Nn \l__zhvt_shukou_page_number_start_dim
    { \l__zhvt_shukou_chapter_title_start_dim + \l__zhvt_shukou_chapter_title_height_dim }
}
\ExplSyntaxOff

\begin{document}

% ★★ 版心（书口）排布覆写：把书名放到鱼尾之下（2026-10-04）
%
%   cls 默认版心自上而下：书名（title 静态框，在鱼尾之上）→ 鱼尾 → 篇名（chapter 动态框）→ 页码
%   本模板改为：          鱼尾（位置不动）→ 书名 → 篇名 → 页码分割线 → 页码
%
%   手法：把 cls 放在 title 静态框里的书名，改放进 chapter 动态框（鱼尾之下）；
%         title 静态框留空。书名取 \l__zhvt_title_str，即 \title{...} 的第一参数，
%         所以换书时只改 \title，不必在此重复写书名。
%
%   ⚠ 三条必须遵守的写法（2026-10-04 实测）：
%     ① 汉字必须用 \hbox{...} 包住。flowfram 会把框内容装进一个宽 14.4mm（= 一个 line_sep
%        = 书口栏宽）的纵向 minipage 并按该宽度折行，而整页靠 PDF /Rotate 90 显示、
%        框内文字本身不变向 —— 不包 \hbox 时「围炉夜话」4 字会被逐字折成 4 行，
%        在版面上就成了横排 4 字并横向溢出版框。
%     ② 竖排下**一个 \hbox 就是一整列**。书名与篇名要同处一列，必须放进同一个 \hbox；
%        分成两个 \hbox 时，后一个会被当成新的一列、甩到隔壁栏去
%        （实测「序」跑到书口栏外侧 14.4mm 的正文列上）。中间的 \quad 在框内是横向空白，
%        显示出来正好是书名与篇名之间的纵向间隔。
%     ③ ★ 内容盒要写成 **与框等宽的定宽盒** `\hbox to \l__zhvt_shukou_chapter_title_height_dim{…\hfil}`。
%        实测：chapter 框是把内容盒沿纵向**居中**摆放的，内容宽度一变，书名就跟着漂：
%        自然宽度的 \hbox（书名 38.4mm）落在 x=157.2mm，加篇名后（57.6mm）落在 x=152.4mm，
%        目录页（无篇名）与正文页的书名就差半格。做成等宽盒后，无论有没有篇名，
%        内容固定从框顶起排，实测两页均落在 x=142.8mm，不再漂。
\ExplSyntaxOn
% 把书口 chapter 框（书名+篇名）起始位置下移 1 字高，避免 4 字卷名顶到鱼尾
% ⚠ cls 的 \__zhvt_calculate_dimensions: 在 \AtBeginDocument 里执行，会重算该 dim；
%   必须把覆写挂到 begin document 之后，否则会被二次覆盖（2026-10-05 实测）。
% (旧方案已移到导言区，直接修改 chapter 框起始位置)
% 卷标题命令 #1=书口篇名(鱼尾) #2=正文卷名（含卷序，卷名内空格用 ~）
%
% ★★ 目录页两项修订（2026-10-05，愚大人指出）★★
%   ① 页码起点统一对齐：
%      原状 —— 目录项写成 \par\hspace{4em}卷名\hspace{4em}页码，页码起点 = 卷名实宽 + 4em。
%      卷名实宽随字数变：「卷一~梁惠王上」6 字 +~、「卷七~离娄上」5 字 +~，
%      差 1 个字位 → 卷七~卷十 4 条的页码整体比其余 10 条高 27.21pt（一个字），看着上下不齐。
%      改法 —— 卷名补全角空格（宽 = 1 字 = 27.215pt）补足到 **7 字符位**再写进目录，
%      14 条页码起点遂全部落在同一条竖线上（起点对齐；末字随页码字数自然参差）。
%      ⚠ 只对目录里的副本 \l_tmpc_tl 补，正文卷名仍用 \l_tmpa_tl 原样 —— 正文像素不受影响。
%   ② 卷名整体上移半字位：前导 \hspace 由 4em 改 3.5em。
%      实测原状目录块在版框内偏下（天头留白 139.6pt / 地脚留白 110.7pt），
%      上移 0.5em = 13.6pt 后上下基本均分，视觉居中。
\NewDocumentCommand { \juan } { mm } {
	\clearpage
	\tl_set:Nn \l__zhvt_chapter_title_tl { #1 }
	\tl_set:Nn \l_tmpa_tl { #2 }
	\tl_set_eq:NN \l_tmpc_tl \l_tmpa_tl
	\int_step_inline:nn { 7 - \str_count:n { #2 } }
	  { \tl_put_right:Nn \l_tmpc_tl { 　 } }
	\tl_set:Nx \l_tmpb_tl { \zhnum { page } }
	\tl_put_left:Nn \l_tmpb_tl { \hspace { 4em } }
	\tl_put_left:Nn \l_tmpb_tl { \l_tmpc_tl }
	\tl_put_left:Nn \l_tmpb_tl { \par \hspace { 3.5em } }
	\addtocontents { toc } { \l_tmpb_tl }
	\skip_horizontal:n { \l__zhvt_font_size_dim * 4 }
	\l_tmpa_tl
	\par
	\thispagestyle { empty }
}
\cs_set_protected:Npn \__zhvt_set_shukou_content: {
	% ★ 鱼尾上方留白（title 静态框）嵌入闲叟钤印：
	%   整页 /Rotate 90，故此处把图用 \rotatebox{90} 转正为竖章；
	%   尺寸卡在版心条幅内（书内约 7mm 宽 × 33.8mm 高，留余量不压鱼尾）。
	\setstaticcontents * { title } {
		\vbox_to_ht:nn { \l__zhvt_baseline_skip_dim } {
			\vskip 14.2pt
			\hbox_to_wd:nn { \l__zhvt_shukou_title_height_dim } {
				\hfil
				\rotatebox [ origin = c ] { 90 } { \includegraphics [ width = 7mm , height = 33.8mm ] { yshs.png } }
				\hfil
			}
		}
	}
% ★ 版心（书口）统一改法（2026-11-16 愚大人定）：18pt＋框内居中＋书口文字同版框色
%   ① 原 27.2pt 时 7 字长篇名（如「廉颇蔺相如列传」）末字溢出 chapter 框、压住页码；
%      改 \fontsize{18pt}{18pt}，前后 \hfil 令字串在框内居中（不随篇名字数漂移）。
%   ② \vbox{\vskip 8.2pt} 定量下移，抵消缩字号后栏宽方向 8.2pt 上偏（像素实测 y 中心 543.3 = 原版）。
%      ⚠ \raisebox / \vspace* 在 flowfram 框内无效；换字号须同步换补偿：16pt→10.1、19pt→7.4。
%   ③ chapter 与 page~number 均置 \color{sealbrown}（HTML 874434），与版框／界行同色；正文仍黑。
%   ⚠ 页码字位、钤印、鱼尾、界行、正文一律未动；五部合订本同此口径。
	\setdynamiccontents * { chapter }
	{ \color{sealbrown} \vbox { \vskip 8.2pt \hbox to \l__zhvt_shukou_chapter_title_height_dim
		{ \hfil \fontsize{18pt}{18pt}\selectfont \l__zhvt_title_str \quad \l__zhvt_chapter_title_tl \hfil } } }
	\setdynamiccontents * { page~number } { \color{sealbrown} \quad \zhnum { page } }
}
\__zhvt_set_shukou_content:
\ExplSyntaxOff

\maketitle

% 单篇，无多卷目录，不渲染目录页

\juan{田叔列传}{卷一百零四~田叔列传第四十四}
\hspace*{2em}田叔者。赵陉城人也。其先。齐田氏苗裔也。叔喜剑。学黄老术于乐巨公所。叔为人刻廉自喜。喜游诸公。赵人举之赵相赵午。午言之赵王张敖所。赵王以为郎中。数岁。切直廉平。赵王贤之。未及迁。\\
\hspace*{2em}会陈豨反代。汉七年。高祖往诛之。过赵。赵王张敖自持案进食。礼恭甚。高祖箕踞骂之。是时赵相赵午等数十人皆怒。谓赵王曰。﹁王事上礼备矣。今遇王如是。臣等请为乱。﹂赵王齧指出血。曰。﹁先人失国。微陛下。臣等当蟲出。公等柰何言若是。毋复出口矣。﹂于是贯高等曰。﹁王长者。不倍德。﹂卒私相与谋弑上。会事发觉。汉下诏捕赵王及群臣反者。于是赵午等皆自杀。唯贯高就系。是时汉下诏书。﹁赵有敢随王者罪三族。﹂唯孟舒。田叔等十馀人赭衣自髡钳。称王家奴。随赵王敖至长安。贯高事明白。赵王敖得出。废为宣平侯。乃进言田叔等十馀人。上尽召见。与语。汉廷臣毋能出其右者。上说。尽拜为郡守。诸侯相。叔为汉中守十馀年。会高后崩。诸吕作乱。大臣诛之。立孝文帝。\\
\hspace*{2em}孝文帝既立。召田叔问之曰。﹁公知天下长者乎。﹂对曰。﹁臣何足以知之。﹂上曰。﹁公。长者也。宜知之。﹂叔顿首曰。﹁故云中守孟舒。长者也。﹂是时孟舒坐虏大入塞盗劫。云中尤甚。免。上曰。﹁先帝置孟舒云中十馀年矣。虏曾一人。孟舒不能坚守。毋故士卒战死者数百人。长者固杀人乎。公何以言孟舒为长者也。﹂叔叩头对曰。﹁是乃孟舒所以为长者也。夫贯高等谋反。上下明诏。赵有敢随张王。罪三族。然孟舒自髡钳。随张王敖之所在。欲以身死之。岂自知为云中守哉。汉与楚相距。士卒罢敝。匈奴冒顿新服北夷。来为边害。孟舒知士卒罢敝。不忍出言。士争临城死敌。如子为父。弟为兄。以故死者数百人。孟舒岂故驱战之哉。是乃孟舒所以为长者也。﹂于是上曰。﹁贤哉孟舒。﹂复召孟舒以为云中守。\\
\hspace*{2em}后数岁。叔坐法失官。梁孝王使人杀故吴相袁盎。景帝召田叔案梁。具得其事。还报。景帝曰。﹁梁有之乎。﹂叔对曰。﹁死罪。有之。﹂上曰。﹁其事安在。﹂田叔曰。﹁上毋以梁事为也。﹂上曰。﹁何也。﹂曰。﹁今梁王不伏诛。是汉法不行也。如其伏法。而太后食不甘味。卧不安席。此忧在陛下也。﹂景帝大贤之。以为鲁相。\\
\hspace*{2em}鲁相初到。民自言相。讼王取其财物百馀人。田叔取其渠率二十人。各笞五十。馀各搏二十。怒之曰。﹁王非若主邪。何自敢言若主。﹂鲁王闻之大惭。发中府钱。使相偿之。相曰。﹁王自夺之。使相偿之。是王为恶而相为善也。相毋与偿之。﹂于是王乃尽偿之。\\
\hspace*{2em}鲁王好猎。相常从入苑中。王辄休相就馆舍。相出。常暴坐待王苑外。王数使人请相休。终不休。曰。﹁我王暴露苑中。我独何为就舍。﹂鲁王以故不大出游。\\
\hspace*{2em}数年。叔以官卒。鲁以百金祠。少子仁不受也。曰。﹁不以百金伤先人名。﹂\\
\hspace*{2em}仁以壮健为卫将军舍人。数从击匈奴。卫将军进言仁。仁为郎中。数岁。为二千石丞相长史。失官。其后使刺举三河。上东巡。仁奏事有辞。上说。拜为京辅都尉。月馀。上迁拜为司直。数岁。坐太子事。时左相自将兵。令司直田仁主闭守城门。坐纵太子。下吏诛死。仁发兵。长陵令车千秋上变仁。仁族死。陉城今在中山国。\\
\hspace*{2em}太史公曰。孔子称曰﹁居是国必闻其政﹂。田叔之谓乎。义不忘贤。明主之美以救过。仁与余善。余故并论之。\\
\hspace*{2em}褚先生曰。臣为郎时。闻之曰田仁故与任安相善。任安。荥阳人也。少孤贫困。为人将车之长安。留。求事为小吏。未有因缘也。因占著名数。武功。扶风西界小邑也。谷口蜀刬道近山。安以为武功小邑。无豪。易高也。安留。代人为求盗亭父。后为亭长。邑中人民俱出猎。任安常为人分麋鹿雉兔。部署老小当壮剧易处。众人皆喜。曰。﹁无怂也。任少卿分别平。有智略。﹂明日复合会。会者数百人。任少卿曰。﹁某子甲何为不来乎。﹂诸人皆怪其见之疾也。其后除为三老。举为亲民。出为三百石长。治民。坐上行出游共帐不办。斥免。\\
\hspace*{2em}乃为卫将军舍人。与田仁会。俱为舍人。居门下。同心相爱。此二人家贫。无钱用以事将军家监。家监使养恶齧马。两人同床卧。仁窃言曰。﹁不知人哉家监也。﹂任安曰。﹁将军尚不知人。何乃家监也。﹂卫将军从此两人过平阳主。主家令两人与骑奴同席而食。此二子拔刀列断席别坐。主家皆怪而恶之。莫敢呵。\\
\hspace*{2em}其后有诏募择卫将军舍人以为郎。将军取舍人中富给者。令具鞍马绛衣玉具剑。欲入奏之。会贤大夫少府赵禹来过卫将军。将军呼所举舍人以示赵禹。赵禹以次问之。十馀人无一人习事有智略者。赵禹曰。﹁吾闻之。将门之下必有将类。传曰﹃不知其君视其所使。不知其子视其所友﹄。今有诏举将军舍人者。欲以观将军而能得贤者文武之士也。今徒取富人子上之。又无智略。如木偶人衣之绮绣耳。将柰之何。﹂于是赵禹悉召卫将军舍人百馀人。以次问之。得田仁。任安。曰。﹁独此两人可耳。馀无可用者。﹂卫将军见此两人贫。意不平。赵禹去。谓两人曰。﹁各自具鞍马新绛衣。﹂两人对曰。﹁家贫无用具也。﹂将军怒曰。﹁今两君家自为贫。何为出此言。鞅鞅如有移德于我者。何也。﹂将军不得已。上籍以闻。有诏召见卫将军舍人。此二人前见。诏问能略相推第也。田仁对曰。﹁提桴鼓立军门。使士大夫乐死战斗。仁不及任安。﹂任安对曰。﹁夫决嫌疑。定是非。辩治官。使百姓无怨心。安不及仁也。﹂武帝大笑曰。﹁善。﹂使任安护北军。使田仁护边田穀于河上。此两人立名天下。\\
\hspace*{2em}其后用任安为益州刺史。以田仁为丞相长史。\\
\hspace*{2em}田仁上书言。﹁天下郡太守多为奸利。三河尤甚。臣请先刺举三河。三河太守皆内倚中贵人。与三公有亲属。无所畏惮。宜先正三河以警天下奸吏。﹂是时河南。河内太守皆御史大夫杜父兄子弟也。河东太守石丞相子孙也。是时石氏九人为二千石。方盛贵。田仁数上书言之。杜大夫及石氏使人谢。谓田少卿曰。﹁吾非敢有语言也。原少卿无相诬汙也。﹂仁已刺三河。三河太守皆下吏诛死。仁还奏事。武帝说。以仁为能不畏彊御。拜仁为丞相司直。威振天下。\\
\hspace*{2em}其后逢太子有兵事。丞相自将兵。使司直主城门。司直以为太子骨肉之亲。父子之间不甚欲近。去之诸陵过。是时武帝在甘泉。使御史大夫暴君下责丞相﹁何为纵太子﹂。丞相对言﹁使司直部守城门而开太子﹂。上书以闻。请捕系司直。司直下吏。诛死。\\
\hspace*{2em}是时任安为北军使者护军。太子立车北军南门外。召任安。与节令发兵。安拜受节。入。闭门不出。武帝闻之。以为任安为详邪。不傅事。何也。任安笞辱北军钱官小吏。小吏上书言之。以为受太子节。言﹁幸与我其鲜好者﹂。书上闻。武帝曰。﹁是老吏也。见兵事起。欲坐观成败。见胜者欲合从之。有两心。安有当死之罪甚众。吾常活之。今怀诈。有不忠之心。﹂下安吏。诛死。\\
\hspace*{2em}夫月满则亏。物盛则衰。天地之常也。知进而不知退。久乘富贵。祸积为祟。故范蠡之去越。辞不受官位。名传后世。万岁不忘。岂可及哉。后进者慎戒之。\\

% ===== 索隐述赞（司马贞《史记索隐》）· 统一样式 =====
% 题头「索隐述赞曰」五字并入反白框；正文每四句一列，定宽盒 \hfil 居列中；
% \fontsize{24}{41} 列距对齐全书界行（121+41k）；\clearpage 使述赞独占一页。
\clearpage
\shuzanhead{索隐述赞曰}\linebreak
{\shuzanbody\fontsize{24}{41}\selectfont
\hbox to \hsize{\hfil 田叔长者。重义轻生。张王既雪。汉中是荣。\hfil}\linebreak
\hbox to \hsize{\hfil 孟舒见废。抗说相明。按梁以礼。相鲁得情。\hfil}\linebreak
\hbox to \hsize{\hfil 子仁坐事。刺举有声。\hfil}%
\par}
\end{document}
